\newpage
\section{单变量函数的连续性}
\label{sec:functions-continuity}

函数的连续性把自变量的接近与函数值的接近联系起来。在某一点附近，这种联系可以保证函数值不突然改变符号；在一个区间上，它又能保证函数取得两个函数值之间的所有值。第一章研究数列时建立的收敛性工具，将帮助我们从局部的定义得到这些整体的结论。

\subsection{连续性的定义}
\label{subsec:functions-continuity-definition}

\begin{definition}
设\(f\)在\(x_0\)及其附近有定义。若\(\lim_{x\to x_0}f(x)=f(x_0)\)，称\(f\)在\(x_0\)处{\heiti{}连续}。也就是说，对任意\(\varepsilon>0\)，都存在\(\delta>0\)，使
\begin{equation}\label{eq:functions-continuity-definition}
 |x-x_0|<\delta\quad\Longrightarrow\quad|f(x)-f(x_0)|<\varepsilon.
\end{equation}
\end{definition}

与极限的定义相比，这里已经给出了极限应等于的数\(f(x_0)\)。当\(x=x_0\)时，上式自然成立，所以不必排除这一点。写成增量的形式，就是\(\Delta x\to0\)时，\(f(x_0+\Delta x)-f(x_0)\to0\)。

定义中的\(\varepsilon\)是函数值的误差要求，\(\delta\)是为此选定的自变量范围，可以随\(\varepsilon\)和\(x_0\)改变。证明时，通常先估计函数值之差，再反过来确定这个范围；选定的\(\delta\)要对范围内所有自变量同时有效。这样的\(\delta\)并不唯一，取更小的正数也都适用。

若只从\(x_0\)左侧或右侧接近，便分别得到{\heiti{}左极限}与{\heiti{}右极限}，记为
\begin{equation*}
 \lim_{x\to x_0^-}f(x),\qquad \lim_{x\to x_0^+}f(x).
\end{equation*}
相应地，只在\(x<x_0\)或\(x>x_0\)的范围内要求式\eqref{eq:functions-continuity-definition}，就是左连续或右连续。对区间内点，函数连续当且仅当左右极限都等于该点的函数值。

若函数在区间的每一点都连续，称它在该区间上连续。闭区间\([a,b]\)的端点分别按右连续与左连续理解，考察的自变量始终在区间内。例如，\(\sqrt x\)在\([0,+\infty)\)上连续，在\(x=0\)处只须从右侧考察。

\begin{example}
设
\begin{equation*}
 f(x)=\begin{cases}
 x+1,&x<0,\\
 a,&x=0,\\
 x^2+b,&x>0.
 \end{cases}
\end{equation*}
求\(f\)在\(x=0\)处的左右极限，并分别求它在该点左连续、右连续及连续的条件。改变\(a\)能否改变左右极限？
\end{example}

判断连续性时，应分别考察点附近的变化与该点的取值。两侧极限相同而函数值不同，可以通过改变一个点的取值恢复连续；两侧极限不同，则不能这样处理。

\begin{exercise}
\mbox{}\par\nobreak

（1）用定义证明\(f(x)=x^2\)在\(x=1\)处连续。可先限制\(|x-1|<1\)，再控制\(|x+1|\)。

（2）当\(x\ne1\)时，设\(g(x)=(x^2-1)/(x-1)\)。应怎样定义\(g(1)\)，才能使它在\(1\)处连续？

\end{exercise}

\subsection{连续函数的基本性质}
\label{subsec:functions-continuity-properties}

第一章的数列极限运算法则，是通过误差估计得到的。同样的估计也适用于函数极限，因此得到以下结论。

\begin{theorem}\label{thm:functions-continuity-operations}
若\(f,g\)在\(x_0\)处连续，则它们的和、差、积也在\(x_0\)处连续；当\(g(x_0)\ne0\)时，商\(f/g\)也在该点连续。若\(g\)在\(x_0\)处连续，\(f\)在\(g(x_0)\)处连续，则复合函数\(f\circ g\)在\(x_0\)处连续。
\end{theorem}

加法、乘法等可以沿用第一章的证明。对于复合，先确定外层允许的误差，再由内层保证它。给定\(\varepsilon>0\)，先对外层函数选取\(\eta>0\)，使
\begin{equation*}
 |u-g(x_0)|<\eta\quad\Longrightarrow\quad
 |f(u)-f(g(x_0))|<\varepsilon.
\end{equation*}
再对内层函数选取\(\delta>0\)，使\(|x-x_0|<\delta\)时有\(|g(x)-g(x_0)|<\eta\)。把\(u=g(x)\)代入，就得到所需结论。

常值函数与\(f(x)=x\)由定义知连续，所以多项式函数连续，有理函数在分母不为零的地方连续。这样，前面分段函数两侧的极限都可以由相应表达式在连接点的取值求得。

\begin{proposition}[局部保号性]\label{prop:functions-continuity-sign}
若\(f\)在\(x_0\)处连续且\(f(x_0)>0\)，则在\(x_0\)的某个邻域内都有\(f(x)>0\)。负值情形类似。
\end{proposition}

\begin{proof}
在连续性定义中取\(\varepsilon=f(x_0)/2\)，便有\(f(x)>f(x_0)/2>0\)。
\end{proof}

同样，连续性还保证函数在一点附近有界：取\(\varepsilon=1\)，便有\(|f(x)|<|f(x_0)|+1\)。这一结论只在局部成立，例如\(f(x)=x\)在每点连续，却在\(\mathbb R\)上无界。

这说明连续函数在非零值附近保持符号。将它用于\(f-g\)，也可知：若\(f(x_0)>g(x_0)\)，则在充分小的邻域内仍有\(f(x)>g(x)\)。

\begin{proposition}\label{prop:functions-continuity-sequence}
若\(f\)在\(x_0\)处连续，且定义域中的数列\(x_n\to x_0\)，则\(f(x_n)\to f(x_0)\)。
\end{proposition}

\begin{proof}
对给定误差，连续性提供\(\delta\)；由\(x_n\to x_0\)，充分靠后的项都满足\(|x_n-x_0|<\delta\)，因而相应的函数值也满足误差要求。
\end{proof}

第一章在递推式中取极限时，已经使用过这种关系。这里也可以把它理解为：连续函数将趋向同一点的自变量数列，变为趋向相应函数值的数列。

\begin{example}
设\(g\)在\(x=0\)附近有界。定义\(f(0)=0\)，\(x\ne0\)时令\(f(x)=xg(x)\)。用连续性定义证明\(f\)在\(x=0\)处连续。是否需要\(g\)本身在\(x=0\)处连续？
\end{example}

\subsection{介值定理}
\label{subsec:functions-intermediate-value}

在区间上连续，意味着函数从一个值变化到另一个值时，不会跳过中间的值。这个几何认识可以由第一章的二分法与区间套得到证明。

\begin{theorem}[介值定理]\label{thm:functions-intermediate-value}
设\(f\)在\([a,b]\)上连续。若\(k\)介于\(f(a)\)与\(f(b)\)之间，则存在\(\xi\in[a,b]\)，使\(f(\xi)=k\)。当\(k\)严格介于两端值之间时，可取\(\xi\in(a,b)\)。
\end{theorem}

\begin{proof}
端点值等于\(k\)时结论已成立，以下设\(f(a)<k<f(b)\)。令\(g(x)=f(x)-k\)。从\([a,b]\)开始二分，若中点处\(g=0\)，便找到所求的点；否则保留左端为负、右端为正的一半。这样得到闭区间套\([a_n,b_n]\)，其长度为\((b-a)/2^n\)，且\(g(a_n)<0<g(b_n)\)。

由闭区间套定理，\(a_n,b_n\)趋于同一点\(\xi\)。连续性给出\(g(a_n),g(b_n)\to g(\xi)\)，故\(g(\xi)\le0\)且\(g(\xi)\ge0\)，于是\(g(\xi)=0\)。两端值次序相反时同理。
\end{proof}

取\(k=0\)，立即得到常用的零点存在定理。

\begin{corollary}[零点存在定理]\label{cor:functions-zero-existence}
若\(f\)在\([a,b]\)上连续，且\(f(a)f(b)<0\)，则\(f\)在\((a,b)\)内至少有一个零点。若它还严格单调，则该零点唯一。
\end{corollary}

端点异号给出存在性，严格单调给出唯一性。二分法还可以逐步缩小零点所在区间；进行\(n\)次二分后，取区间中点近似零点，误差不超过\((b-a)/2^{n+1}\)。

\begin{example}
证明方程\(x^3+x-1=0\)在\((0,1)\)内有唯一实根。不使用导数，怎样通过两点函数值之差判断严格单调？再写出进行三次二分后保留的区间。
\end{example}

\begin{example}
设\(f\)在\([0,1]\)上连续，且对所有\(x\in[0,1]\)都有\(0\le f(x)\le1\)。证明存在\(\xi\in[0,1]\)，使\(f(\xi)=\xi\)。这个结论是否保证从任意初值出发的迭代都收敛？
\end{example}

若区间上的连续函数处处不为零，它就只能始终为正或始终为负：否则，由异号的两点就能得到零点。这将局部保号性推进为整个区间上的判断。

\begin{exercise}
\mbox{}\par\nobreak

（1）证明对任意实数\(t\)，方程\(x^3+x=t\)都有唯一实根。端点不再固定时，应怎样根据\(t\)选择一个有根的区间？

（2）设\(f,g\)在\([a,b]\)上连续，且\(f(a)<g(a)\)、\(f(b)>g(b)\)。证明两函数的图像在\((a,b)\)内相交。

（3）端点同号的连续函数能否在区间内部有零点？若没有连续性，端点异号是否仍能保证存在零点？分别举例说明。

（4）（2013年重庆卷理科第6题改编）若\(a<b<c\)，函数\(f(x)=(x-a)(x-b)+(x-b)(x-c)+(x-c)(x-a)\)的两个零点分别位于哪两个相邻区间内？
\end{exercise}

\begingroup
\setlength{\parskip}{4pt}
\subsection{闭区间上的连续函数}
\label{subsec:functions-extreme-values}

介值定理说明，连续函数不会遗漏两个函数值之间的值。接下来考虑取值范围的两端：函数值是否有界？如果有界，是否一定存在最大的函数值和最小的函数值？

\begin{example}
比较\(f(x)=x\)在\((0,1)\)与\([0,1]\)上的取值，再考察\(g(x)=1/x\)在\((0,1]\)上的取值。它们分别是否有界，是否有最大值和最小值？
\end{example}

有界与取得最值是两件事。例如，\(x\in(0,1)\)时，\(f(x)=x\)的值始终小于\(1\)，却没有最大的函数值；每选定一个函数值，总还能找到更大的。若把两个端点都包含进来，最大值与最小值就都能取到。下面的定理表明，闭区间上连续函数的一般情形也是如此。

\begin{theorem}[最值定理]\label{thm:functions-extreme-values}
若\(f\)在闭区间\([a,b]\)上连续，则它在该区间上有界，且存在\(u,v\in[a,b]\)，使对所有\(x\in[a,b]\)，都有
\begin{equation*}
 f(u)\le f(x)\le f(v).
\end{equation*}
也就是说，\(f\)在\([a,b]\)上有最小值和最大值。
\end{theorem}

函数值可以越来越接近某个高度，但这个高度是否对应图像上的一点？证明的关键是把函数值的逼近与自变量的收敛联系起来：先选取逼近这个高度的函数值，再从对应的自变量中找出收敛子列，最后利用连续性确定极限处的函数值。

\begin{proof}
先证有界。若函数无界，便能选出\(x_n\in[a,b]\)，使\(|f(x_n)|>n\)。这些自变量未必收敛，但它们都在同一个闭区间内，故可取收敛子列\(x_{n_j}\to c\in[a,b]\)。连续性给出\(f(x_{n_j})\to f(c)\)，与\(|f(x_{n_j})|>n_j\)矛盾。

有界之后，还须证明函数值能够到达最高位置。记\(M\)为全部函数值的最小上界，即\(M\)是上界，而任何更小的数都不是上界。这样的数可由二分得到：在一个函数值与某个上界之间取中点，若中点仍是上界，就向下缩小范围；否则向上缩小，闭区间套便确定\(M\)。

现在先选函数值逼近\(M\)，再寻找对应自变量的聚集位置。由最小上界的意义，可取\(y_n\in[a,b]\)，使\(M-1/n<f(y_n)\le M\)。抽取收敛子列\(y_{n_j}\to v\in[a,b]\)，连续性使函数值也趋于\(f(v)\)，而夹逼原理给出同一极限为\(M\)，故\(f(v)=M\)。对\(-f\)使用此结论，即得最小值。
\end{proof}

这里用到了第一章的结论：闭区间中的自变量数列有收敛子列，且极限仍在区间内；连续性则把自变量的极限与函数值的极限联系起来。两者共同保证了逼近的高度能够真正取到。

若最小值为\(m\)、最大值为\(M\)，再由介值定理，函数的值域恰为\([m,M]\)。所以要判断\(f(x)=k\)是否在闭区间内有解，可以转化为判断\(m\le k\le M\)。

\paragraph{恒成立与存在性} \mbox{}

设\(f\)在\([a,b]\)上连续，最小值与最大值分别为\(m,M\)。则
\begin{equation*}
 \begin{aligned}
 \text{对所有 }x\in[a,b],\ f(x)\ge t
 &\quad\Longleftrightarrow\quad t\le m,\\
 \text{存在 }x\in[a,b],\ f(x)\ge t
 &\quad\Longleftrightarrow\quad t\le M.
 \end{aligned}
\end{equation*}
前者要求控制全部取值，后者只须有一个取值满足条件。对于严格不等式，最值能够取到也必须计入判断。

\begin{example}
设\(f(x)=x^2-2x+2\)，\(x\in[0,3]\)。用配方法求值域，并分别求实数\(t\)的范围，使下列条件成立：

（1）对所有\(x\in[0,3]\)，都有\(f(x)\ge t\)。

（2）存在\(x\in[0,3]\)，使\(f(x)<t\)。

（3）方程\(f(x)=t\)在\([0,3]\)上有解，并进一步讨论不同\(t\)对应的根的个数。
\end{example}

\begin{example}
求实数\(t\)的范围，使不等式\(x^2-2tx+2\ge0\)对所有\(x\in[0,2]\)成立。再求存在\(x\in[0,2]\)使该表达式小于零的\(t\)的范围。将参数分离与直接求最小值两种方法进行比较。
\end{example}

\begin{exercise}
\mbox{}\par\nobreak

（1）求\(h(x)=x^2-4x+5\)在\([0,3]\)上的值域，并分别求使\(h(x)>t\)恒成立、至少存在一个\(x\)使\(h(x)>t\)的参数范围。

（2）设\(f\)在\([a,b]\)上连续且处处为正。证明存在常数\(c>0\)，使所有函数值都不小于\(c\)。将闭区间换为开区间，结论是否仍成立？

（3）设\(f\)在\([0,1]\)上连续，且\(f(0)=f(1)\)。证明存在\(\xi\in[0,1/2]\)，使\(f(\xi+1/2)=f(\xi)\)。应对哪个函数使用零点存在定理？
\end{exercise}

\begin{example}
设连续函数\(f,g\)在两个闭区间\(I,J\)上的值域分别为\([m,M]\)、\([p,P]\)。分别求实数\(t\)的范围，使下列条件成立：

（1）对所有\(x\in I\)、\(y\in J\)，都有\(f(x)\ge t+g(y)\)。

（2）对每个\(x\in I\)，都存在\(y\in J\)，使\(f(x)\ge t+g(y)\)。

（3）存在\(x\in I\)，使对所有\(y\in J\)，都有\(f(x)\ge t+g(y)\)。
\end{example}

判断这些条件时，应分别确定两个自变量各自受哪个量词约束，再用相应的最大值或最小值表达条件。

\begin{exercise}
\mbox{}\par\nobreak

（1）（2012年福建卷文科第15题）已知关于\(x\)的不等式\(x^2-ax+2a>0\)在\(\mathbb R\)上恒成立，则实数\(a\)的取值范围是\underline{\hspace{1.5cm}}。

（2）（2025年天津卷第15题）已知\(a,b\in\mathbb R\)，且\((2a+b)x^2+bx-a-1\le0\)对任意\(x\in[-2,2]\)成立，则\(2a+b\)的最小值为\underline{\hspace{1.2cm}}。
\end{exercise}

\endgroup
\begingroup
\setlength{\parskip}{3pt}
\subsection{连续函数的反函数}
\label{subsec:functions-continuous-inverse}

讨论反函数时，先要判断能否由函数值唯一确定自变量；在此基础上，再考察这种反过来的对应是否连续。对于区间上的连续函数，这两个问题可以联系起来研究。

\begin{theorem}\label{thm:inv-mono}
设\(f\)在区间\(I\)上连续。它在\(I\)上存在反函数的充要条件是严格单调。当此条件成立时，\(f^{-1}\)在其定义域\(f(I)\)上也连续，且与\(f\)具有相同的单调方向。
\end{theorem}

\begin{proof}
严格单调保证不同自变量对应不同的值，所以反函数存在。反过来，若反函数存在，连续图像就不能发生折返，否则会重复经过某个高度。具体地，对\(x_1<x_2<x_3\)，若\(f(x_2)>\max\{f(x_1),f(x_3)\}\)，取介于这个最大值与\(f(x_2)\)之间的\(k\)，介值定理便分别在两段内给出函数值为\(k\)的点，与一一对应矛盾。中间值低于两端的情况同理。因此中间点的函数值总严格介于两端值之间，固定一对点的增减方向后，再插点或向两侧延伸都保持该方向，故\(f\)严格单调。

再证明反函数连续。反函数以\(y\)为自变量、以\(x\)为函数值，所以先限制\(x\)的误差，再寻找能保证它的\(y\)的范围。设\(f\)严格递增，固定\(y_0=f(x_0)\)。在\(x_0\)两侧取\(x_-,x_+\in I\)，使它们与\(x_0\)的距离小于给定的\(\varepsilon\)。由单调性，\(f(x_-)<y_0<f(x_+)\)。取\(\delta\)不大于两侧函数值与\(y_0\)的距离，则
\begin{equation*}
 |y-y_0|<\delta\quad\Longrightarrow\quad
 x_-<f^{-1}(y)<x_+\quad\Longrightarrow\quad
 |f^{-1}(y)-x_0|<\varepsilon.
\end{equation*}
端点处只保留相应一侧的估计，递减情形同理。由介值定理，\(f(I)\)也是区间，以上估计均在反函数的定义域内使用。
\end{proof}

研究反函数不一定要先写出它的表达式。若\(f\)递增，由原函数值的大小就能估计反函数值：
\begin{equation*}
 f(a)<y<f(b)\quad\Longleftrightarrow\quad a<f^{-1}(y)<b.
\end{equation*}
求\(f^{-1}(y)\)就是求方程\(f(x)=y\)的唯一解，可以用二分法逐步确定；反函数连续则说明解随\(y\)连续变化。这样，方程的存在性、唯一性以及解随右端的变化，就能在同一个函数关系中得到解释。

\begin{example}
证明\(f(x)=x^3+x\)在\(\mathbb R\)上有连续反函数，求其反函数的定义域和值域。若不求出反函数的表达式，能否判断\(\lim_{y\to2}f^{-1}(y)\)？
\end{example}

\begin{exercise}
\mbox{}\par\nobreak

（1）将\(f(x)=x^2\)分别限制在\([-2,0]\)与\([0,2]\)上，写出两种反函数，并说明它们的单调方向与连续性。

（2）构造一个区间上的严格单调函数，使其值域不是区间。它的图像怎样说明，严格单调本身不能代替连续性？

\end{exercise}

\begin{exercise}
\mbox{}\par\nobreak

（1）（2015年上海卷文科第4题）设\(f^{-1}(x)\)为\(f(x)=x/(2x+1)\)的反函数，则\(f^{-1}(2)=\underline{\hspace{1.2cm}}\)。

（2）（2013年上海卷理科第14题）对区间\(I\)上有定义的函数\(g(x)\)，记\(g(I)=\{y\mid y=g(x),x\in I\}\)。已知定义域为\([0,3]\)的函数\(y=f(x)\)有反函数\(y=f^{-1}(x)\)，且\(f^{-1}([0,1))=[1,2)\)，\(f^{-1}((2,4])=[0,1)\)。若方程\(f(x)-x=0\)有解\(x_0\)，则\(x_0=\underline{\hspace{1.2cm}}\)。
\end{exercise}

\endgroup
\newpage
\subsection{连续性的数列刻画*}
\label{subsec:functions-sequential-continuity}

前面已知，连续函数保持数列的极限。反过来，若所有趋近于同一点的自变量数列都给出相应的函数值极限，也足以确定连续性。

\begin{theorem}\label{thm:functions-sequential-continuity}
函数\(f\)在定义域中的点\(x_0\)处连续，当且仅当：对定义域内任意满足\(x_n\to x_0\)的数列，都有\(f(x_n)\to f(x_0)\)。
\end{theorem}

\begin{proof}
必要性已由命题\ref{prop:functions-continuity-sequence}证明。若\(f\)不连续，则存在\(\varepsilon_0>0\)，使无论把邻域取得多小，其中总能找到一点，其函数值与\(f(x_0)\)的距离不小于\(\varepsilon_0\)。对每个\(n\)，取这样的一点\(x_n\)，使
\begin{equation*}
 |x_n-x_0|<\frac1n,\qquad
 |f(x_n)-f(x_0)|\ge\varepsilon_0.
\end{equation*}
于是\(x_n\to x_0\)，函数值却不趋于\(f(x_0)\)，与条件矛盾。
\end{proof}

这里把定义中“任意小的邻域”转化为一列逐步缩小的邻域。它既能用来证明连续，也能通过适当选择趋近数列来否定连续；后者尤其适合在同一点附近有不同变化方式的函数。

使用这个刻画证明连续时，必须考察任意趋于\(x_0\)的数列；仅沿某一列点得到正确极限，并不能控制其他趋近方式。否定连续性则只须找到一列点，使相应函数值不趋于\(f(x_0)\)。上面的证明还说明，不连续时总能让自变量越来越接近\(x_0\)，而函数值始终与\(f(x_0)\)相隔至少一个固定的正数。

\begin{example}
设\(f(x)=\sin(1/x)\)，\(x\ne0\)。分别选取趋于零的数列，使函数值恒等于\(1\)与\(-1\)，由此说明无论怎样定义\(f(0)\)，都不能使它在\(x=0\)处连续。再比较\(g(x)=x\sin(1/x)\)、\(g(0)=0\)的情形。
\end{example}

\begin{exercise}
设\(f\)在\([0,1]\)上连续，数列\(x_n\)满足\(x_{n+1}=f(x_n)\)，且各项都在\([0,1]\)内。

（1）若\(x_n\)收敛，证明其极限必为\(f\)的不动点。

（2）取\(f(x)=1-x\)，讨论不同初值产生的数列，并说明不动点的存在为什么不能保证迭代收敛。

（3）若偶数项与奇数项分别趋于\(u,v\)，证明\(f(u)=v\)、\(f(v)=u\)。此时整个数列收敛的充要条件是什么？结合上一问解释\(u\ne v\)时的行为。
\end{exercise}

\subsection{一致连续*}
\label{subsec:functions-uniform-continuity}

在一点处讨论连续性时，先固定基点，再按函数值的误差要求选取自变量的范围。改变基点，允许的范围也可能改变。若要同时控制整个区间内的变化，就需要进一步问：能否选取同一个范围，使它对所有基点都适用？

\begin{definition}
设函数\(f\)定义在\(A\subseteq\mathbb R\)上。若对任意\(\varepsilon>0\)，都存在\(\delta>0\)，使对任意\(x,y\in A\)，都有
\begin{equation*}
 |x-y|<\delta\quad\Longrightarrow\quad|f(x)-f(y)|<\varepsilon,
\end{equation*}
则称\(f\)在\(A\)上{\heiti{}一致连续}。
\end{definition}

这里的\(\delta\)可以随\(\varepsilon\)改变，但必须对\(A\)中任意一对点同时有效。一致连续自然保证每一点处连续；反过来，分别为每一点找到合适的\(\delta\)，却未必能得到一个对所有点都有效的正数。

例如，对\(f(x)=x^2\)，当\(x,y\in[0,1]\)时，
\begin{equation*}
 |f(x)-f(y)|=|x-y||x+y|\le2|x-y|.
\end{equation*}
给定\(\varepsilon>0\)，取\(\delta=\varepsilon/2\)，就能同时保证整个区间内的误差要求。因此\(x^2\)在\([0,1]\)上一致连续。

但在\(\mathbb R\)上，取\(x_n=n\)、\(y_n=n+1/n\)，便有
\begin{equation*}
 |x_n-y_n|=\frac1n\to0,\qquad
 |f(x_n)-f(y_n)|=2+\frac1{n^2}>2.
\end{equation*}
所以，对\(\varepsilon=1\)，无论选取多小的\(\delta>0\)，总能找到距离小于\(\delta\)而函数值相差大于\(1\)的两点。这个函数在每一点都连续，却在\(\mathbb R\)上不一致连续。

同一个函数在不同范围上表现不同。对于闭区间，前面用于证明最值定理的收敛子列工具，还能把逐点的连续性推进为整个区间上的一致连续性。

\begin{theorem}\label{thm:functions-uniform-continuity}
闭区间上的连续函数一致连续。
\end{theorem}

如果结论不成立，就能找到越来越接近、函数值却始终相隔一定距离的两列点。闭区间保证它们有共同的聚集位置，连续性则使两列函数值必须趋向同一个数，由此得到矛盾。

\newpage
\begin{proof}
设\(f\)在\([a,b]\)上连续，却不一致连续。则存在\(\varepsilon_0>0\)，使对每个正整数\(n\)，都能选出\(x_n,y_n\in[a,b]\)，满足
\begin{equation*}
 |x_n-y_n|<\frac1n,\qquad |f(x_n)-f(y_n)|\ge\varepsilon_0.
\end{equation*}
由第一章的收敛子列定理，可取\(x_{n_j}\to c\in[a,b]\)。又因为
\begin{equation*}
 |y_{n_j}-c|\le|y_{n_j}-x_{n_j}|+|x_{n_j}-c|\to0,
\end{equation*}
所以\(y_{n_j}\to c\)。由连续性，\(f(x_{n_j})\)与\(f(y_{n_j})\)都趋于\(f(c)\)，故它们之差趋于零，与上述不等式矛盾。
\end{proof}

这一结论使我们能够把闭区间分成足够短的小段，同时控制每一段内函数值的波动。以后用有限求和研究整体累积时，这种统一的控制将成为估计总误差的依据。

一致连续也可以用数列刻画。前一小节固定一个趋近位置，研究函数值是否趋于该位置的函数值；这里则允许两列点的位置不断移动，只要求它们彼此越来越接近。

\begin{proposition}\label{prop:functions-uniform-continuity-sequences}
函数\(f\)在\(A\)上一致连续，当且仅当：对\(A\)中任意两列满足\(|x_n-y_n|\to0\)的点，都有\(|f(x_n)-f(y_n)|\to0\)。
\end{proposition}

\begin{proof}
若\(f\)一致连续，给定\(\varepsilon>0\)，选取定义中的\(\delta\)。由\(|x_n-y_n|\to0\)，充分靠后的项都满足\(|x_n-y_n|<\delta\)，因而函数值之差小于\(\varepsilon\)。

反过来，若\(f\)不一致连续，就能像前面定理的证明一样，选出\(x_n,y_n\in A\)，使\(|x_n-y_n|<1/n\)，而\(|f(x_n)-f(y_n)|\ge\varepsilon_0>0\)，与条件矛盾。
\end{proof}

这里不要求两列点各自收敛，例如\(n\)与\(n+1/n\)都趋于无穷，却彼此接近。闭区间定理的证明则利用收敛子列，将移动的比较归到同一点，再使用该点的连续性。

\begin{exercise}
\mbox{}\par\nobreak

（1）证明\(f(x)=\sqrt x\)在\([0,+\infty)\)上一致连续。可先证明\(|\sqrt x-\sqrt y|\le\sqrt{|x-y|}\)。

（2）判断\(g(x)=1/x\)在\([1,+\infty)\)与\((0,1)\)上是否一致连续，并给出证明。

（3）设\(f\)在\(A\)上一致连续，\(x_n\)是\(A\)中的Cauchy数列。证明\(f(x_n)\)也是Cauchy数列。用\(g(x)=1/x\)及\(x_n=1/(n+1)\)说明，逐点连续不能保证这一结论。
\end{exercise}

